% Сквозная нумерация внутри раздела (лекции): Теорема 1.1, Определение 1.2, ...
% refname задаёт, как окружение называется в \cref{...}.
\declaretheorem[name=Теорема, numberwithin=section,
  refname={теорема,теоремы}, Refname={Теорема,Теоремы}]{theorem}
\declaretheorem[name=Лемма, sibling=theorem,
  refname={лемма,леммы}, Refname={Лемма,Леммы}]{lemma}
\declaretheorem[name=Утверждение, sibling=theorem,
  refname={утверждение,утверждения}, Refname={Утверждение,Утверждения}]{proposition}
\declaretheorem[name=Следствие, sibling=theorem,
  refname={следствие,следствия}, Refname={Следствие,Следствия}]{corollary}

\declaretheorem[name=Определение, style=definition, sibling=theorem,
  refname={определение,определения}, Refname={Определение,Определения}]{definition}
\declaretheorem[name=Пример, style=definition, sibling=theorem,
  refname={пример,примеры}, Refname={Пример,Примеры}]{example}
\declaretheorem[name=Упражнение, style=definition, sibling=theorem,
  refname={упражнение,упражнения}, Refname={Упражнение,Упражнения}]{exercise}

\declaretheorem[name=Замечание, style=remark, numbered=no]{remark}
\declaretheorem[name=Примечание, style=remark, numbered=no]{note}

\renewcommand{\proofname}{Доказательство}

% Прочие названия для \cref
\crefformat{equation}{(#2#1#3)}
\crefname{section}{лекция}{лекции}
\crefname{figure}{рис.}{рис.}

% Рамка для ключевых формул/фактов
\newtcolorbox{keybox}[1][]{
  colback=blue!3, colframe=blue!50!black, boxrule=0.6pt, arc=2pt,
  fonttitle=\bfseries, title={#1}
}
